% Generated by scripts/build_reports.py; edit reviewed Markdown inputs.
% Compile with XeLaTeX.
% Options for packages loaded elsewhere
\PassOptionsToPackage{unicode}{hyperref}
\PassOptionsToPackage{hyphens}{url}
\PassOptionsToPackage{dvipsnames,svgnames,x11names}{xcolor}
\documentclass[
  11pt,
]{article}
\usepackage{xcolor}
\usepackage[margin=1in]{geometry}
\usepackage{amsmath,amssymb}
\setcounter{secnumdepth}{-\maxdimen} % remove section numbering
\usepackage{iftex}
\ifPDFTeX
  \usepackage[T1]{fontenc}
  \usepackage[utf8]{inputenc}
  \usepackage{textcomp} % provide euro and other symbols
\else % if luatex or xetex
  \usepackage{unicode-math} % this also loads fontspec
  \defaultfontfeatures{Scale=MatchLowercase}
  \defaultfontfeatures[\rmfamily]{Ligatures=TeX,Scale=1}
\fi
\usepackage{lmodern}
\ifPDFTeX\else
  % xetex/luatex font selection
\fi
% Use upquote if available, for straight quotes in verbatim environments
\IfFileExists{upquote.sty}{\usepackage{upquote}}{}
\IfFileExists{microtype.sty}{% use microtype if available
  \usepackage[]{microtype}
  \UseMicrotypeSet[protrusion]{basicmath} % disable protrusion for tt fonts
}{}
\makeatletter
\@ifundefined{KOMAClassName}{% if non-KOMA class
  \IfFileExists{parskip.sty}{%
    \usepackage{parskip}
  }{% else
    \setlength{\parindent}{0pt}
    \setlength{\parskip}{6pt plus 2pt minus 1pt}}
}{% if KOMA class
  \KOMAoptions{parskip=half}}
\makeatother
\usepackage{longtable,booktabs,array}
\usepackage{caption}
\captionsetup[table]{skip=6pt}
\newcounter{none} % for unnumbered tables
\usepackage{calc} % for calculating minipage widths
% Correct order of tables after \paragraph or \subparagraph
\usepackage{etoolbox}
\makeatletter
\patchcmd\longtable{\par}{\if@noskipsec\mbox{}\fi\par}{}{}
\makeatother
% Allow footnotes in longtable head/foot
\IfFileExists{footnotehyper.sty}{\usepackage{footnotehyper}}{\usepackage{footnote}}
\makesavenoteenv{longtable}
\setlength{\emergencystretch}{3em} % prevent overfull lines
\providecommand{\tightlist}{%
  \setlength{\itemsep}{0pt}\setlength{\parskip}{0pt}}
\usepackage{amsmath,amssymb}
\usepackage{xurl}
\usepackage{fancyhdr}
\pagestyle{fancy}
\fancyhf{}
\fancyhead[L]{\small Pure Thought | Revised research specifications}
\fancyfoot[L]{\small 2026-09-17}
\fancyfoot[R]{\small\thepage}
\setlength{\headheight}{15pt}
\setlength{\emergencystretch}{3em}
\allowdisplaybreaks
\usepackage{bookmark}
\IfFileExists{xurl.sty}{\usepackage{xurl}}{} % add URL line breaks if available
\urlstyle{same}
% fallback for those not using the hyperref driver hyperxmp:
\makeatletter
\@ifundefined{xmpquote}{\newcommand{\xmpquote}[1]{#1}}{}
\makeatother
\hypersetup{
  pdftitle={27. Central configurations in a bounded{,} symmetry-reduced N-body problem},
  pdfauthor={Pure Thought research portfolio},
  colorlinks=true,
  linkcolor={blue},
  filecolor={Maroon},
  citecolor={Blue},
  urlcolor={blue},
  pdfcreator={LaTeX via pandoc}}

\title{27. Central configurations in a bounded, symmetry-reduced N-body
problem}
\author{Pure Thought research portfolio}
\date{Revised 17 September 2026}

\begin{document}
\maketitle
\thispagestyle{empty}
\clearpage

\section{27. Central configurations in a bounded, symmetry-reduced
N-body problem}

\textbf{Revised specification --- 2026-09-11.}
\href{https://github.com/rick-stevens-ai/Pure-Thought/blob/main/reviews/27-N-Body-Central-Configurations.md}{Detailed
critique} \ensuremath{\cdot}
\href{https://github.com/rick-stevens-ai/Pure-Thought/blob/main/archive/original-PRDs/27-N-Body-Central-Configurations.md}{Original
description} \ensuremath{\cdot}
\href{https://github.com/rick-stevens-ai/Pure-Thought/blob/main/reviews/EVIDENCE-STANDARD.md}{Evidence
standard}

\subsection{Plain-language guide}

\subsubsection{The problem in everyday terms}

Most collections of gravitating bodies change their overall shape as
they move. Certain special arrangements can rotate or expand and
contract while keeping the same shape. These are central configurations.
This problem asks for every such arrangement in a carefully limited
case, with a proof that the list is complete and a separate analysis of
what small disturbances do to the associated rotating motion.

\subsubsection{Key terms}

\begin{itemize}
\item
  \textbf{N-body problem} --- The problem of describing the motion of N
  mutually gravitating bodies.
\item
  \textbf{Central configuration} --- An arrangement in which each body's
  gravitational acceleration is proportional to its position relative to
  the center of mass, with one common proportionality factor.
\item
  \textbf{Homographic motion} --- Motion preserving the arrangement's
  shape while its size and orientation may change.
\item
  \textbf{Relative equilibrium} --- A motion that is stationary in an
  appropriate rotating reference frame.
\item
  \textbf{Center of mass} --- The mass-weighted average position, used
  here to remove overall translation.
\item
  \textbf{Normalization} --- A fixed choice of scale, such as setting
  the moment of inertia to one, so scaled copies are not counted
  separately.
\item
  \textbf{Symmetry quotient} --- The convention for treating
  configurations related by rotations, reflections or specified
  relabelings as equivalent.
\item
  \textbf{Collision-free} --- Having strictly positive separation
  between every pair of bodies.
\item
  \textbf{Coriolis term} --- A velocity-dependent term arising in a
  rotating frame, essential to the stability calculation.
\item
  \textbf{Spectral stability} --- A statement about eigenvalues of the
  linearized dynamics; it is weaker than a full nonlinear stability
  proof.
\end{itemize}

\subsubsection{Why this matters}

Central configurations are organizing structures in gravitational
dynamics and provide exact reference cases amid otherwise complicated
motion. A complete small catalog can test numerical methods and reveal
how mass choices and symmetry affect possible motions. The problem also
links gravity to exact algebraic root finding.

\subsubsection{What progress would mean}

Progress would establish a complete list under a stated mass choice and
symmetry convention. It would not classify every periodic orbit: a
choreography such as the figure-eight changes shape and belongs to a
different question.

\subsection{Core question}

For fixed positive masses and a declared symmetry quotient, can all
collision-free central configurations in the chosen family be certified
and their relative-equilibrium linearization analyzed?

\subsection{Scope and assumptions}

Start with three labeled positive rational masses in the plane. Fix G=1,
center of mass zero, moment of inertia
I=\ensuremath{\Sigma}m\_i|r\_i|\ensuremath{{}^{2}}=1, and quotient by
O(2), so mirror images count together. Permute only equal masses when
explicitly requested. Extend to a specified four-body symmetry family
after validating the complete three-body case.

Central configurations generate homographic motions; general periodic
choreographies are a separate problem.

\subsection{Mathematical target}

Use positive
U=\ensuremath{\Sigma}\_\{i<j\}m\_im\_j/|r\_i\ensuremath{-}r\_j| and

\[\nabla_{r_i}U=-\lambda m_ir_i,\qquad \lambda =U/I>0.\]

Introduce distance variables d\_ij\textgreater0 with
d\_ij\ensuremath{{}^{2}}=|r\_i\ensuremath{-}r\_j|\ensuremath{{}^{2}},
and retain positivity/collision exclusions when polynomializing. Use
symmetry charts covering every configuration claimed. Completeness
requires certified global domain coverage or an exact
elimination/root-count argument, not just isolated numerical solutions.

\subsection{Required outputs and proof obligations}

Return normalized configurations, distance enclosures, residual
identities, existence/uniqueness proofs, equivalence representatives,
and completeness scope. For planar relative equilibria analyze the
reduced rotating-frame first-order system including Coriolis terms.
Distinguish spectral stability from nonlinear stability.

\subsection{Validation and rejection controls}

For three labeled masses under the stated O(2) quotient, recover three
collinear orderings (one per middle body) and one equilateral shape.
Check \ensuremath{\lambda}=U at I=1. Test collision artifacts introduced
by denominator clearing. For stability, reproduce a known three-body
criterion using identical mass and reduction conventions.

\subsection{Milestones and research extension}

Baseline: complete certified three-body catalog. Strong: complete
enumeration in one four-body symmetry family or mass instance, with all
charts covered. Research extension: a five-body restricted family; no
claim to solve the general finiteness problem without an appropriately
global proof.

Milestones are acceptance gates, not calendar promises. Before a run,
freeze the concrete instance/domain, tolerance, compute budget, and
permitted dependencies. A bounded search that fails to find a result
must preserve an unresolved outcome.

\subsection{Starting instruction}

Work on the baseline above. First make every input and convention
explicit, then validate the reference case. Build the checker alongside
the calculation. Label numerical evidence, finite-model certificates,
and continuum theorems separately; return the strongest justified result
and identify the exact unmet proof obligation.

\subsection{Primary references and their role}

\begin{itemize}
\tightlist
\item
  \href{https://arxiv.org/abs/2103.08906}{Yu and Zhu, On the finiteness
  of four-body central configurations} --- Context for finiteness versus
  explicit certified enumeration.
\end{itemize}

References support definitions or limitations; the scope choices and
proposed milestones are editorial recommendations, not claims that these
results have already been obtained.

\subsection{Displayed formulation and claim boundary}

\href{https://github.com/rick-stevens-ai/Pure-Thought/blob/main/PRDs/27-N-Body-Central-Configurations.md}{Canonical
specification}.

\[U=\sum_{i<j}\frac{m_im_j}{|r_i-r_j|},\quad
\sum_i m_i r_i=0,\quad I=\sum_i m_i|r_i|^2=1,\]

\[\nabla_{r_i}U=-\lambda m_i r_i,\qquad\lambda=U/I>0,
\qquad d_{ij}^2=|r_i-r_j|^2,\quad d_{ij}>0.\]

Preserve collision exclusions when clearing denominators. For three
labeled positive masses modulo \(O(2)\), recover three collinear
orderings and one equilateral shape. Complete coverage requires more
than isolated root boxes. Relative-equilibrium stability uses the
reduced rotating-frame system including Coriolis terms; a figure-eight
choreography is not a central configuration.

\subsection{Detailed review and rationale}

\href{https://github.com/rick-stevens-ai/Pure-Thought/blob/main/PRDs/27-N-Body-Central-Configurations.md}{Revised
problem} \ensuremath{\cdot}
\href{https://github.com/rick-stevens-ai/Pure-Thought/blob/main/archive/original-PRDs/27-N-Body-Central-Configurations.md}{Preserved
original} \ensuremath{\cdot}
\href{https://github.com/rick-stevens-ai/Pure-Thought/blob/f854d4e12b55c1f5b5c2a964b75f2d23496d981e/PRDs/27-N-Body-Central-Configurations.md}{Pinned
upstream source}

Line references below refer to the preserved original at commit
\texttt{f854d4e12b55}, not the rewritten file. ``Blocking'' means the
original issue can invalidate the scientific target or claimed
certificate; ``major'' means it materially changes scope, inference, or
acceptance criteria. These are critiques of specifications, not results
of running the embedded research-code sketches.

\subsubsection{Detailed findings}

\begin{enumerate}
\def\labelenumi{\arabic{enumi}.}
\item
  \textbf{Blocking --- wrong example (26--29, 120).} The figure-eight
  choreography changes shape and is not a central configuration. It
  belongs in a periodic-orbit problem, retained separately in the
  supplement.
\item
  \textbf{Blocking --- counts and quotients (112--114).} The text
  confuses five restricted-problem Lagrange points with Euler central
  configurations. Counting depends on labels, rotations and reflections;
  the rewrite states the quotient and a four-shape three-body benchmark.
\item
  \textbf{Blocking --- stability shortcut (124--130).} Hessian signs
  alone do not give rotating-frame orbital stability; Coriolis terms and
  symmetry modes must be treated.
\item
  \textbf{Major --- algebraic and finiteness claims (36, 93--98).}
  Clearing denominators can add collisions and does not remove square
  roots without auxiliary variables. Numerical roots do not prove
  completeness, and the generic-five-body status claim needs
  qualification against the literature. The new bounded scope avoids an
  unsupported universal claim.
\end{enumerate}

\subsubsection{What the revision preserves and changes}

The central motivation is retained, but the executable question is now:
\textbf{For fixed positive masses and a declared symmetry quotient, can
all collision-free central configurations in the chosen family be
certified and their relative-equilibrium linearization analyzed?} The
revised scope chooses a concrete starting model; those choices are
proposed research-design decisions, not facts inferred from the
original. The baseline can succeed without resolving the full open
research program.

The replacement also removes the original implementation sketches from
the active specification. They remain in the archive for comparison, but
have not been made executable or certified. New implementations should
derive their equations from the revised definitions rather than copy
unchecked prototypes.

\subsubsection{Acceptance criteria to use instead}

Baseline: complete certified three-body catalog. Strong: complete
enumeration in one four-body symmetry family or mass instance, with all
charts covered. Research extension: a five-body restricted family; no
claim to solve the general finiteness problem without an appropriately
global proof.

\subsubsection{Sources supporting the corrections}

\begin{itemize}
\tightlist
\item
  \href{https://arxiv.org/abs/2103.08906}{Yu and Zhu, On the finiteness
  of four-body central configurations} --- Context for finiteness versus
  explicit certified enumeration.
\end{itemize}

\section{Evidence and execution standard}

This standard applies to the revised specifications. It is a proposed
contract for future research runs; this editorial revision does not
supply the scientific certificates themselves.

\subsection{What ``pure thought'' means here}

Use explicitly stated mathematical models and deductions without fitting
conclusions to experimental outcomes. Synthetic examples, exact model
inputs, analytic references, numerical search, and established software
are allowed. Physical constants, electronic-structure approximations,
constitutive laws, and empirical parameter sets must be identified when
used. A model-derived result is conditional on that model.

The original ban on legacy software conflicts with its proposed reliance
on NumPy, SciPy, CVXPY, nauty, and other libraries. Reimplementing every
optimizer and algebra system is not necessary for scientific
independence. Prefer a small independently checkable certificate and a
documented trusted base. A fresh implementation can be a project
objective, but must not be confused with a mathematical proof or novelty
claim.

\subsection{Four evidence levels}

{\def\LTcaptype{none} % do not increment counter
\begin{longtable}[]{@{}
  >{\raggedright\arraybackslash}p{(\linewidth - 4\tabcolsep) * \real{0.3333}}
  >{\raggedright\arraybackslash}p{(\linewidth - 4\tabcolsep) * \real{0.3333}}
  >{\raggedright\arraybackslash}p{(\linewidth - 4\tabcolsep) * \real{0.3333}}@{}}
\toprule\noalign{}
\begin{minipage}[b]{\linewidth}\raggedright
Label
\end{minipage} & \begin{minipage}[b]{\linewidth}\raggedright
What it establishes
\end{minipage} & \begin{minipage}[b]{\linewidth}\raggedright
What it does not establish
\end{minipage} \\
\midrule\noalign{}
\endhead
\bottomrule\noalign{}
\endlastfoot
Numerical evidence & A reproducible computation at stated precision,
sample size, and discretization & Exact truth, continuum coverage, or a
universal theorem \\
Finite-model certificate & An exact/validated statement about the
specified finite model or finite relaxation & An omitted tail,
thermodynamic limit, or physical realization \\
Conditional mathematical theorem & A proof over the stated domain with
every hypothesis and error controlled & Truth outside those assumptions
or a demonstrated physical system \\
Formal proof & A checked theorem in a named proof assistant, with
imported axioms/dependencies listed & Correct physical modeling merely
by virtue of formalization \\
\end{longtable}
}

Formalization is an optional extra deliverable unless the specific
project makes it mandatory. High precision, JSON output, solver status
and small residuals do not by themselves confer any of these proof
levels.

\subsection{Claim directions}

For minimization over a larger feasible set, the optimum is a lower
bound on the original minimum. A physical feasible point supplies an
upper bound. For maximization the directions reverse. Track feasibility
of the primal/dual witnesses separately from their objective values.

Infeasibility of an inner ansatz excludes only that ansatz.
Infeasibility of an outer relaxation can exclude the original problem if
the map into the relaxation is proved and all approximations preserve
necessity. Passing necessary tests means ``not excluded''; existence
needs a constructive sufficient argument.

A smallest observed example is not a minimality theorem. A complete
catalog must define the universe, equivalence relation, boundaries, and
proof that pruning misses nothing. A finite-size numerical trend does
not prove an asymptotic limit.

\subsection{Minimum run manifest}

Before calculation, record the problem ID and revision hash, exact
inputs and conventions, quantified domain, equivalence relation, target
observable, accepted evidence level, tolerances, and resource limits.
Record dependencies and precision/rounding mode. Where randomization is
used, preserve seeds, trial counts and raw sufficient statistics.
Resource limits are planning decisions, not evidence that a mathematical
solution cannot exist.

An output bundle should include:

\begin{itemize}
\tightlist
\item
  The statement being claimed, including all assumptions and
  quantifiers.
\item
  Construction or proof object, plus an independently runnable verifier
  and its trust assumptions.
\item
  Numerical/error ledger separating roundoff, discretization,
  truncation, optimization, sampling and modeling error where relevant.
\item
  Positive reference cases and negative controls that the verifier must
  reject.
\item
  Coverage: included domain, omitted domain, unresolved regions and
  termination reason.
\end{itemize}

These are requirements for future implementations, not a promise of
polynomial-size certificates for every problem.

\subsection{Checking specific certificate types}

\textbf{LP/SDP/SOS:} verify equality constraints, objective direction,
and cone feasibility using exact arithmetic or outward-rounded
intervals. Rationalizing a numerical solution can break PSD constraints.
Correct residuals with justified norm bounds; do not silently project
and reuse the old objective. Distinguish a proof of a polynomial
approximation from a proof for the original function.

\textbf{Interval numerics:} state function enclosures, domains and
outward rounding. A small root box needs existence and, if claimed,
uniqueness. Cover the complement for a global enumeration. Include
singularities and tails rather than discarding them numerically.

\textbf{SAT/SMT:} a solver's UNSAT label is not a certificate. Preserve
a proof supported by a checker and the proof that the encoding, symmetry
reduction and blocked solutions match the mathematical universe. DRAT
applies to an appropriate Boolean CNF proof, not automatically to every
real-arithmetic or SDP claim.

\textbf{Statistics:} specify the estimand, dependence assumptions,
confidence procedure and bias. Zero observed failures does not prove
zero failure probability. MCMC error bars based on estimated
autocorrelation are not generally rigorous mixing guarantees.

\subsection{Milestones and novelty}

Baseline means a working, independently validated result on a concrete
reference problem. Strong means a broader or sharper result with its
claimed evidence level. Research extension means a possible novel result
whose feasibility and significance need review. A negative theorem can
be valuable; an unsuccessful search can also be useful if its scope is
documented, but the two are not interchangeable.

Do not promise ``publication-quality'' in a fixed number of months.
Before claiming novelty, conduct a focused current literature comparison
against the exact theorem/model/metric. The sources in this revision
check key definitions and limitations; they are not an exhaustive
novelty audit across all 30 fields.

\end{document}
